\begin{exercisechunk}
\item \textbf{Density does not determine an approximation rate.}
\label[exercise]{ex:l15-density-versus-rate}
\exerciseprofile{2}{2}{2}{\exproof\quad\excalculation\quad\exinterpretation}{%
Distinguish density of an approximation family from the rate for an
individual target.}
Let $\cH=\ell_2(\mathbb N)$ with standard orthonormal basis
$(e_j)_{j\geq1}$, and let
$V_m=\operatorname{span}\{e_1,\ldots,e_m\}$.
\begin{enumerate}[(a),beginpenalty=10000]
\item Prove that $\bigcup_{m\geq1}V_m$ is dense in $\cH$.
\item Fix $p>0$ and define $f^{(p)}\in\cH$ by
$f_j^{(p)}=j^{-p-1/2}$. Prove that
\[
 \inf_{v\in V_m}\|f^{(p)}-v\|_2
 =\left(\sum_{j>m}j^{-2p-1}\right)^{1/2}
 \asymp m^{-p}
 \qquad\text{as }m\to\infty.
\]
\item Define $f^{(\log)}_1=0$ and
\[
 f^{(\log)}_j=\frac1{\sqrt j\log j},\qquad j\geq2.
\]
Prove that $f^{(\log)}\in\cH$ and
\[
 \inf_{v\in V_m}\|f^{(\log)}-v\|_2
 \asymp\frac1{\sqrt{\log m}}
 \qquad\text{as }m\to\infty.
\]
\item Explain which statement is common to all targets in parts (b) and (c),
and which quantitative conclusion depends on the target. Use the two examples
to state precisely why density alone does not imply a polynomial approximation
rate.
\end{enumerate}
\begin{hints}
The best approximation in $V_m$ is the orthogonal projection onto the first
$m$ coordinates. Compare each tail sum with the corresponding improper
integral.
\end{hints}
\end{exercisechunk}
